% =====================================================================
%  R. Nielsen, SPIRITUAL REALITY, I-II (Spiritual Reality; More on Spiritual Reality)
%  In: For Idea and Reality.  A Journal, edited by R. Nielsen, Bj. Bjørnson and Rud. Schmidt.  Year 1870. Second Volume (Bind 3) and Year 1871. Second Volume (Bind 5).
%  Copenhagen: Chr. Steen & Søn, Publishers; Printed by Louis Klein -- Otto B. Wroblewsky, Publisher; Printed by G. S. Wibe.  Bind 3, pp. 289-320, Bind 5, pp. 1-38.  (Aandsvirkelighed; Videre om Aandsvirkelighed, in For Ide og Virkelighed.)
%  ENGLISH TRANSLATION, from transcription.tex (Danish), the text of record.
%  GENERATED by ebook-method/for-idee-1870-73/tr_build.py --write from this template and the translators' parts in
%  ebook-method/for-idee-1870-73/tr/ (TRANSLATOR-BRIEF.md; terms: GLOSSARY.md, its key terms approved by Hans
%  2026-10-13/14).  Conventions: TRANSLATION-PLAYBOOK.md.  Page numbers of the print in the margin.
% =====================================================================
\documentclass[12pt,a4paper,openany]{book}
\usepackage[utf8]{inputenc}
\usepackage[T1]{fontenc}
\usepackage{libertinus}
\usepackage{libertinust1math}
\usepackage{textalpha}
\usepackage[english]{babel}
\usepackage[protrusion=true,expansion=true]{microtype}
\usepackage{setspace}
\usepackage[margin=1.3in]{geometry}
\usepackage{fancyhdr}
\usepackage{hyperref}
\hypersetup{pdftitle={Spiritual Reality}, pdfauthor={Rasmus Nielsen}, hidelinks}
\setlength{\headheight}{15pt}
\pagestyle{fancy}
\fancyhf{}
\fancyhead[LE,RO]{\thepage}
\fancyhead[C]{\textit{R. Nielsen: Spiritual Reality (1870-71)}}
\setstretch{1.3}
\usepackage{marginnote}
\renewcommand{\thefootnote}{\arabic{footnote})}   % the print numbers each page's notes 1), 2)
\newcommand{\opage}[1]{\marginnote{\footnotesize\textit{[#1]}}}
\newcommand{\essayhead}[2]{\begin{center}{\large #1}\ifx&#2&\else\\[0.3em]{\normalsize #2}\fi\\[0.5em]\rule{3em}{0.4pt}\end{center}\medskip\noindent}
\newcommand{\signature}[1]{\par\smallskip\hfill #1\hspace*{2em}\par}

\begin{document}

% the title page of Bind 3 (1870-2-for-idee-og-virkelighed.pdf), in English
\begin{titlepage}
\begin{center}
\vspace*{3em}
{\Huge For Idea and Reality.}\\[1em]
\rule{8em}{0.4pt}\\[1em]
{\large A Journal,}\\[0.8em]
{\small edited}\\[0.4em]
{\footnotesize by}\\[0.6em]
{\large R. Nielsen,}\\[0.4em]
{\large Bj. Bjørnson and Rud. Schmidt.}\\[1em]
\rule{8em}{0.4pt}\\[0.6em]
{\large Year 1870. Second Volume.}\\[0.6em]
\rule{8em}{0.4pt}\\[4em]
Copenhagen.\\[0.4em]
Chr. Steen \& Søn, Publishers.\\[1.5em]
{\small Printed by Louis Klein.}\\[3em]
{\small Translated from the Danish}
\end{center}
\end{titlepage}

\essayhead{Spiritual Reality.}{}
% --- p. 289 ---
\setcounter{footnote}{0}%
\opage{289}Revealed religion, with the Word of Life, with miracles
and signs, in historical as well as super-historical facts,
is the work of the spirit. To wish here to reach for
the outward real and forget the spirit, the eternal
inward, is as absurd as to wish to grasp the
pure spirit apart from the facts in which it
has revealed itself. The spirit of life and of truth must
be seen in the light of reality, and reality in turn in
the light of the spirit: the reality of revelation is Spiritual Reality.

As the reality is that is to be seen, so
must the eye be that is to see the real. To
see merely sensible things it is enough to have a sensibly
sound and sharp eye; in this respect the animal
may even have certain one-sided advantages over man. What
especially distinguishes human sight from
animal sight is what is peculiar in man's gaze,
that man looks at the thing thinkingly.

The investigator of nature arms his eye --- with the telescope,
to see the great that would otherwise be hidden in
unattainable distance; with the microscope, to see what,
despite its immediate nearness, lies hidden in
% --- p. 290 ---
\setcounter{footnote}{0}%
\opage{290}infinite smallness. He wants to see thoroughly what there is to see
everywhere in the visible world; but he is not content with
a sensible seeing; sight must serve his thought; he
searches in everything visible for the invisible: he inquires
after laws.

Thought is spiritual; a world of real laws is
a spiritual world. We may say with full right that
the investigator of nature sees the spiritual in the real, that
reality stands before his gaze as a spiritual reality;
but we cannot say that this spiritual
reality is yet Spiritual Reality. And why
not? It is not the fault of reality; no, it is the fault
of the eye or, seen from another side, it is the fault of the illumination.
The spiritual light in which the investigator
alone can, will, and ought to see the whole real realm of nature
is the light of thought, of reason, of science; but this
light of thought, for all its clarity, is too general, cold, and
one-sided to be the light of the spirit.

That in the fruitful regions of reality there is still
something else to discover than what can be seen by the general
light of thought of science, art and
poetry bear witness enough. The landscape painter, for
example, can never grow weary of contemplating nature, for
the longer he contemplates it, the more he sees in
it; it communicates itself to him, it entrusts him with its
secrets: he constantly makes new discoveries.
When the beholder looks at the painted landscape and finds
the picture wonderfully beautiful, he is easily tempted to
say to himself: this is no faithful rendering of
the real; it is nature, not as it really
is, but as it ought to be in the name of the idea and of beauty
% --- p. 291 ---
\setcounter{footnote}{0}%
\opage{291}be; art must re-create reality and transform
nature.

There is a peculiar mixture of true
and false in such judgments. For it is quite true that the work of art
cannot be a straightforward copy of material reality,
the reality that anyone who likes can
go and contemplate with his own eyes; but what is false
is that the artist should have added something of his
own, in the sense that, for the embellishment of nature,
he has put in something that does not belong to the revelation
of nature's essence, something not-natural, and thus
something in reality unnatural. It is so far from being the case
that the nature that has found its transfiguration in art
should be a falsified nature at odds with reality,
that on the contrary it is the innermost essence of nature, the
true nature, that receives its appearance in the true
work of art. The beholder must see the work of art in the light
in which the artist has seen nature. If he has no eye
for the idea of art, then he sees in the work of art an arbitrary
piece of human handiwork, a deceptive imitation, perhaps
even a fraud, and holds fast, sensibly, to the
sensible reality. What good does it do to show
him a Thorvaldsen's Jason! He sees a naked man
with a helmet on his head and a sheepskin on his arm,
and involuntarily asks: did the real Jason
look like that? Did Jason really exist, and did
Thorvaldsen see him? Since one cannot give him an
affirmative answer to either of these questions, he
declares the whole thing a play of the imagination without reality,
and so without truth.

What contributes especially in such inquiries
% --- p. 292 ---
\setcounter{footnote}{0}%
\opage{292}to confusing sound sense is an unclear mixture
of the two concepts: \emph{Reality and Truth.} That
the real must have its truth, and the true always
in one way or another correspond to the real, is
indisputable; but for that very reason it must be insisted
that truth itself stands infinitely higher than the outward,
palpable reality.

If everything real were without more ado truth,
then there could be no real error,
no real crime or real lie: the reality
of error, of lie, of crime would then
be merely different expressions for real truth. To be sure,
one speaks of a true crime, a true lie;
but who does not see that the adjective ``true'' in such connections
precisely stamps the crime as crime,
the lie as lie! Conversely, if there were no truth
except in what exists in an outward palpable way,
then there would be no truth either in thought
or reason, either in soul or spirit. But that
the inner, invisible has its reality just as well, and in
reality its truth, as the outer, palpable,
indeed a far more essential and deeper truth, no
one can deny without denying all truth.

That the truth in the essential inner is of higher
order than the truth in the real outer is precisely
what, on the level of the idea, gives poetry its significance
and divine right. The strongest and boldest
break of poetry with outward reality occurs in the fairy tale. H.
C. Andersen's fairy tale \emph{The Shadow}, for example, has many interesting
features of real life and could almost
take on the appearance of being a real story, if
% --- p. 293 ---
\setcounter{footnote}{0}%
\opage{293}the highly gifted author had not --- it grieves us
to have to tell him so --- allowed himself so highly
offensive a breach of the immutable laws of nature; for
a ``sensible'' reader says (and the same is said by every
boy in an American boys' school): ``Surely no one
is going to make me believe that a shadow can
sneak away from its master like that, creep in under the cake-woman's
skirt, begin to grow fat while the master grows
thin, and, what is the worst scandal and at the same time
the most improbable of all, carry it so far that in the end it gets engaged, think of it --- to a Princess of the
Blood.'' We admit that the author, doubtless out of sheer
ignorance, which can only in part be excused by
his not being a physicist by profession, has overlooked the very
important circumstance that a born shadow, even if it
is the shadow of ``a learned man,'' never, no, however
much it may stretch and extend itself, never
manages to become a real human being, even if it
became a fat human being. If only one could forgive the poet
his ignorance of physics, one could not deny
him excellence through and through in psychology;
for the psychical laws, the well-known natural laws of conceit, of vanity,
of egoism and of humbug,
are fulfilled with the unerring
certainty of true genius most punctiliously in \emph{The Shadow.}

We need not lose ourselves here in lengthy inquiries
concerning the relation between science,
art and poetry; for this much can be clearly
seen from what has been said: 1) that the question of \emph{reality}
must always be subordinated to the question of \emph{truth;}
and 2) that the question of truth every time
% --- p. 294 ---
\setcounter{footnote}{0}%
\opage{294}acquires a peculiar significance in relation to the field of vision
within which the inquiry falls. Truth can in no sphere,
no more in poetry than in science, give up reality;
but truth stands above the outward,
material reality. The most important question is
not: how much reality is there in the true; but,
conversely: how much truth is there in the real?
When for this reason it is insisted again and again that what
is truth in poetry and art is not truth in
science, and conversely: then the meaning is easy to grasp.
Poetry and art cannot and ought not to conceive reality
in the same way as science; the scientific
investigator cannot and ought not to see reality
with the same eyes as the poet. If a zoologist
took it into his head to introduce ducks and geese speaking, in order properly to
bring out these animals' characteristic peculiarities,
one would rightly say: that is an unscientific
procedure; there is no seriousness in such a
treatment. But because the zoologist sins against truth
when he lets the animals speak, who will therefore maintain
that the fairy tale ``\emph{the Ugly Duckling}'' does not have
a deep poetic truth corresponding to reality?
To separate the shadow from its shadow-caster and let it go
its own way is impermissible in optics; but there is
surely no man of cultivation and sound sense
who accuses the poet of violating truth because
he, by a higher plenary authority, breathes life into a shadow-being,
holds a mirror before the reader's eyes and says:
this is how it goes in reality!

It is not merely, as people say, permitted in art
and poetry to alter something in the real and to break
% --- p. 295 ---
\setcounter{footnote}{0}%
\opage{295}with what natural science understands
by the laws of things; but it is an essential requirement, an unavoidable
duty for the poet as for the artist to treat
the material of reality in such a way that there can be
truth in poetry and art. Just as truth in all
circles of the life of consciousness has and must have reality
beneath it, so in all the circles of existence
the laws of truth have the laws of reality beneath them. If
only the laws of truth are fulfilled, then the fulfilment of
the laws of reality follows as of itself; if, conversely, the
laws of outward reality are set above the laws of truth, then
art and poetry are only lies and deceit.

Wherever the question of truth is raised, what is asked for
is a judgment upon reality, and then first of all
the standard by which judgment is to be passed. As regards science,
art, and poetry, the cultivated are on the
whole fairly well oriented; but when the
question concerns the relation between truth and reality
in religion, one might be tempted to believe that we
were still living in the golden age of papacy or at any rate of theology;
for the barbarism is appalling. ``Did Christ
really walk upon the sea? Did he really raise
Lazarus? Did he really ascend into heaven in such an
outward, visible way that I, for example, had I
been present, could have seen it just as plainly
as the apostles? Give me an infallible, historical,
best of all natural-historical proof that it all
really happened as it is told:
then I will believe it. But if no one can give
me full assurance of the reality of the marvel-
% --- p. 296 ---
\setcounter{footnote}{0}%
\opage{296}ous, then I should be a fool to take it
for anything but legend and myth''\footnote[1]{How the outward reality can, in the sacred history, set itself up in a deterrent way against Spiritual Reality, we learn from the Gospel of Christ's walking on the sea. From the Savior's word: \emph{Come}! Peter receives an impression of Spiritual Reality. ``But when he saw the strong wind,'' in other words: when he received an overwhelming impression of the sensible reality of nature, ``he was afraid''\dots{} Mtth. 14:29—30.}.

Those who in our day make such demands upon the
gospel history betray the same lack of cultivation
in respect of religion as those who, in the name of truth,
demand that the picture of Jason shall demonstrably resemble
the real Jason, and the story of the Shadow be a
real story, betray in respect of poetry and art.
Torn loose from the idea of art, reality acquires no,
absolutely no significance in art; torn loose from the spirit
of revelation, reality acquires no, absolutely no significance
in religion.

By this it is by no means said that the real is to be conceived
and judged in religion in the same or even in a
similar way as in poetry and art; far
from it! here precisely the opposite is said; for here it is expressly said
that reality is to be conceived and judged in
another way, by a quite different standard, in
religion than in whatever other of the different
spheres of the life of the spirit, be it in the practical
or the theoretical, be it in art and poetry or in
science.

What the empirically investigating sciences establish
as strict, immutable laws are laws for the outwardly
given reality. Such things, such
% --- p. 297 ---
\setcounter{footnote}{0}%
\opage{297}laws, and conversely: at this conception science
stops; whence the things and the laws have their
origin, it does not investigate. If we raise the question
whether it would not be possible that there could be
another order of things, another nature than this present one,
then science, when it duly reflects,
must answer: it is indeed possible, but then there must also be
other laws. Thereby science's relation to
reality is essentially characterized. Religion may
very well concede the validity of this principle of science;
but that it cannot enter upon
presenting reality in the same illumination as science
lies in the nature of the case. If religion regarded
reality from the same point of view as science,
it would not be religion.

The laws of truth in art and poetry are laws of beauty.
Religion need not take a hostile stand toward the corresponding
conception and treatment of reality;
indeed one may even say that it stands in a
more intimate and familiar relation to poetry than to
scientific investigation of nature. Nonetheless the spirit of
revelation lays quite another weight on real
facts and makes quite differently strict demands of reality
than the aesthetic idea\footnote[1]{The aesthetic ideal belongs solely to the imagination; it satisfies contemplation; but makes no demand upon the will; it requires no realization, it has in illusion its own unreal reality, its completion and conclusion. The religious ideal, by contrast, is in the very strictest sense an ideal of reality, for it is first and last an ideal of the will, and the ideal of the will demands its realization. Philosophy of Religion, p. 205---6.}.

% --- p. 298 ---
\setcounter{footnote}{0}%
\opage{298}That physicists and aestheticians may perhaps find
the category of Spiritual Reality somewhat unclear need not surprise
us; for it is natural that everyone lives himself into
his own and accustoms himself to see reality with his own eye;
but that those who undertake to do the big talking about the scientific
grounding of religion should cast aside the expression
\emph{Spiritual Reality}, and --- as if as a proof
that their own spirit is without reality, and their
own reality without spirit --- try to stamp this
expression as an ambiguity, betrays unambiguously that they
in their burning zeal for theology have forgotten the
first article of faith, the article on creation.

Faith in creation surely presupposes, obviously, that
the whole visible, sensible nature is the work of the spirit.
But if the real nature is originally the work of the spirit,
then it follows indisputably that revealed
religion wants the reality of the whole world regarded
in the light of the spirit, that, in other words, all reality must
from the beginning be conceived as Spiritual Reality. To
avoid ambiguity, attention has therefore also been drawn in its place\footnote[1]{Philosophy of Religion, p. 475---76.}
expressly to the fact that the religious contemplation of nature
is in principle absolutely heterogeneous with the
scientific. ``Without regard to the possible objections of
science, and long before these came into being, revelation
has spoken its word on creation, and ''the religious
consciousness has in the inwardness of faith found inspired
words in praise of the power that upholds all, that
can shake the mountains in their foundations, can roll up the heavens
like a cloth, can draw back its breath, so
% --- p. 299 ---
\setcounter{footnote}{0}%
\opage{299}that all that lives dies\dots{} It looks indeed almost as if
the religious view of nature were just as sensibly objective
as the observation corresponding to science; it looks
as if the believer saw God himself in outward nature,
in the mighty sea, where he builds his lofty halls in the waters,
saw the glance of his eyes in lightning, heard his voice
in thunder, his rolling chariot in storm-clouds, when
he rides upon the wings of the storm. Yet this is not so;
the objectively natural is everywhere only occasional cause,
only means and point of transition for that which is
subjective in the spirit. ''

The spirit of revelation has something to communicate to Adam's
race, something to enlighten men about, and for this it
uses reality. How it uses the
outward real, in what facts, events,
pictures of life and illuminations it makes visible its
power and will, its decrees, promises, and laws,
must therefore be judged solely and alone by the end it has
set itself, the task it wants to solve in religion. In
this respect there are two questions on whose answer
everything depends: 1) \emph{how does the reality
of the events presented in the Bible
stand to the word of Fall and Restoration
in which the spirit reveals the eternal truth?}
2) \emph{to what extent can it agree} with the \emph{infallibility of the spirit
and the credibility of the word that in the
sacred history traits are cited of a view of the world
that manifestly deviates from, and in some
particulars even expressly conflicts with,
the conception now prevailing of real
% --- p. 300 ---
\setcounter{footnote}{0}%
\opage{300}natural conditions, of the order of things
expressed in visible facts?}

Before an honest and thorough reconciliation
between religion and culture can be thought of, it is necessary to bring forward
these questions and see to what extent, under the
conditions of the times, it is possible to get them answered in a way
satisfactory to both parties.

Religion stands above culture; if the race had only
a choice between one of two: either to give up religion and
pay homage to culture, or to hold fast to positive religion
and break with culture, then the more serious among our
contemporaries would surely not hesitate to choose the latter.
We side with those who strictly demand that the revealed
word be done justice in every respect; we place ourselves
on the side of those who, if it were necessary, would rather
break the staff over the whole view of the world founded on natural science
than give up a jot of
that word of faith which manifestly belongs to salvation and blessedness.
It is from this point of view that in this connection we
will make our contribution to the answering of the questions.

Thus: \emph{how does the reality
of the events told in the Bible stand to
the word of revelation on Fall and Restoration?}

To begin at the beginning, we direct
attention to the accounts in the first book of Moses
on creation, Paradise, and the Fall. That these accounts
deal with something real and in the very strictest
sense concern reality, no one will be able to deny.
We say: it is Spiritual Reality, reality for the spirit,
and understand this to mean that the believing conception
has absolutely nothing to do with the outward physical reality
severed from
% --- p. 301 ---
\setcounter{footnote}{0}%
\opage{301}the spirit, but that the believer
who in religious self-understanding wants to understand the word has first and
last to do with the spirit in the word, has to regard
the traits belonging to natural reality as elements,
means, and conditions which the spirit uses to make visible
the truth it wants to reveal.

To establish the correctness of this way of looking at things,
we shall now seek to show that two men can, as regards
the religious content of the narratives, believe the same,
the very same, with the same strength and inwardness
of conviction, and yet have a different conception
of what is outwardly real in the events.

Let us begin with a simple believer. He
takes everything literally; he follows the history of creation
step by step from the first to the last day, and it
never occurs to him that there could be any other
explanation of the real coming-into-being of the world. He reads
the description of the Garden of Eden as one reads a description
of any other real garden, and has
no scruple about the fact that doubt might perhaps arise
concerning the reality of the two supernatural
trees that are here planted in an otherwise natural
soil. He sees Eve before the Tree of Knowledge; he
hears the serpent's tempting words; he becomes a witness to the Fall
and feels in his conscience that the judgment of Jehovah is not
only severe and terrible, but also holy and just;
in short: the believer apprehends in these events
the outward in unity with the inward, the supernatural
in unity with the natural, the real in unity
with the spiritual, hence the whole as Spiritual Reality,
reality for the spirit.

% --- p. 302 ---
\setcounter{footnote}{0}%
\opage{302}If we now place ourselves on the standpoint of cultivation, of the investigation of nature,
and of reasoning criticism, that is, on
the stage of consciousness where there can be talk of a
sharp, intelligent separation of the spiritual and the
real, as of the supernatural and the natural:
then it is absolutely necessary that from the beginning
an account be given of what is decisive for religion and faith,
whether it is the outward real or
the inward, the spiritual; for on the determination of
this point the relation between religion and culture
may be said on the whole to rest.

If the testimonies about the outward real are
to be set higher than the testimony of the spirit, then natural science must
be made supreme judge. ``It can then
easily be proved that the facts cited in the history of creation
can as little be brought into agreement
with what has shown itself to the sharp eye of the investigators
to be reality as the notion of a tree
from whose fruits one should be able to eat one's way to becoming
as wise as God himself can be reconciled with sound
understanding. A spirit that is so unacquainted with the most elementary
principles of natural science
and of sound reason, so completely devoid of astronomical, geological,
botanical, zoological, and anthropological preliminary knowledge,
can impossibly be any real spirit, still
less God's own Spirit; with this the matter is settled.''

But such attacks upon religion, with however much
acuteness and scientific competence they may be
equipped, are without any significance, since they proceed from a
thoroughly false presupposition. In revelation and faith it is
by no means the outward real that stands above the spirit,
% --- p. 303 ---
\setcounter{footnote}{0}%
\opage{303}but conversely it is the spirit itself that is unconditionally lord
over all reality, and therefore uses the real as
a means for its own end. The fundamental question is, for
what end the spirit has uttered its word on creation,
Paradise, and the Fall, what it is that it
wishes to teach us in this its word, and what truth there is in this
teaching.

It is then, in the first place, the reality of creation
that is at issue, but, mark well, conceived as
Spiritual Reality. That the natural, real world
is a revelation of God's almighty, creating word
is the original presupposition of religion and faith, a
presupposition that scientific criticism will surely
leave unshaken. To make this truth comprehensible, graspable,
intelligible to the simple, the various
stages of creation are referred to a creation in six days.

Like the word on creation, so also
the word on Paradise has its reality; but --- it is Spiritual Reality,
reality for the spirit. That man originally
has life from God's life, spirit from God's spirit, and
hence is by his ideal vocation formed to bear
God's image, to know God's commandment and do his will,
this is the supreme truth of anthropology, and real
experience has surely shown that in the same degree as
the consciousness of this is obscured, in the same degree man sinks,
despite cultivation of the understanding and culture, into
animal dependence upon selfish powers of nature, and loses
the impression of the original nobility of his race.

And now the word on temptation and the Fall! Every man,
even the hero who is the object of the admiration of millions,
begins his career with a double con-
% --- p. 304 ---
\setcounter{footnote}{0}%
\opage{304}ception of human life, a conception that it
in reality is what it ought not to be, and a
conception that it in reality ought to be what
it is not; this rift between ideal and reality
the individual can see in race and history only insofar as
he sees it in his own inner being. The hero can now go
two ways. If he believes the word of the lost Paradise,
then he acknowledges himself complicit in the transgression
of the race, then he believes in the reality of the image of God in the
holy prototype and strains with all his powers after the conditions of restoration
and liberation, to work
himself and the race up out of the deep abasement, the false
reality whose essence is bondage, lie, and deceit,
and to win a step forward toward that reality of the spirit
which satisfies the demands of life. If he despises the word
of the lost Paradise, then he chooses another way.
In dim, mythical self-deification he then begins by
making the outward real, with all the glorification
that can be attained thereby, the end of life.
He does not cast contempt outright upon the ideal conceptions;
far from it! The old dreams of the race, the memories of Paradise
of peace and freedom, right and truth, he employs
as a means of awakening for the credulous people, as
exalted, unselfish motives, magic formulas
with which one dazzles the mob. He goes further, he makes
a critical divorce between the real and the ideal.
The real is that he who has the powers uses them
to his own advantage; the ideal is that the world wants to be deceived;
so he smiles at the spirit of the world and says: thy
will be done. Such a heroic figure must necessarily acquire
two faces, an official and a private. The official, which
% --- p. 305 ---
\setcounter{footnote}{0}%
\opage{305}turns outward toward the crowd, has noble features, the stamp of
loftiness and majesty; it is in exalted calm, radiant
with the splendor of the Ideal, a counterfeit, fabricated image of God.
But inside this sits another face, the person's
honest, real private face, turned inward with a satyr's grin
that mocks the one turned outward. In this way
the hero in history comes to resemble the fairy Alcine in
the fairy tale: so long as the enchantment works, and the lovers
are blinded, she looks like a beautiful woman; but
when the sorcery fades, one discovers with horror
that she is an old, sallow, wrinkled hag. Yet even
this mockery of the image of God and of the sacred memories of Paradise
is an eloquent testimony that the word about
temptation, fall and original guilt is reality,
Spiritual Reality, reality for the spirit.

That it is not the representations of the outwardly
real, of the creation in six days, of the trees in
the Garden of Eden, of the appearance of the serpent, and so on, on which
the matter turns, is plainly seen from the fact that these representations
are like lightly cast colored pictures, which may well
support the visualization, but do not in and by themselves
exercise the slightest influence on faith.

Two people can --- we repeat it --- have the
same faith in the divine truth of the Word, and yet
be of widely differing opinion concerning the outward side
of the events. When the one takes everything literally,
while the other grasps the outward features figuratively,
a difference of opinion undeniably shows itself here; but that
this difference of opinion does not concern the Word as a word of revelation,
can be seen from two well-known facts.

First, he who takes everything literally
% --- p. 306 ---
\setcounter{footnote}{0}%
\opage{306}is not thereby in a position to win any new gain
for the life of faith. To believe in the Creation is to believe that this
world has really come into being at the bidding of Almighty power; but
the representation of six days, six weeks, six million
years makes, in a religious respect, neither here nor there. To
believe in Paradise and the Fall is to believe in the truth
of the image of God in which man was originally
created, and, in the pain of self-contradiction, to recognize the
original guilt by which the Ideal was darkened. But
the image of the tree of knowledge and of the speaking serpent can
be set up as objects of faith only through a monstrous misunderstanding.
That the Fall of the innocent presupposes
a temptation is a truth of faith; but whether the tempting
serpent had legs and a human head, whether it
sat in the tree or lay at the root of the tree, what
color it had, whether it was bluish-gray, green
or spotted, is a topic for poetic imagination, an
object of opinion, but no object of faith.

Second, the disagreement between him who takes
everything literally and him who grasps the narrative figuratively
will immediately show itself as a disagreement about
physics and natural history; but precisely the circumstance
that the grounds of inquiry for and against must all be drawn
from the outwardly real, quite as in the inquiries
about the river Phrat, about bdellium and the stone onyx,
is surely a testimony that the contending parties have left
the foundation of revelation and betaken themselves out onto the territory of natural science,
which has no concern with faith.

But in granting the spirit the right to
use reality as it wills, for the illustration
of the truth of its word, we must above all not
% --- p. 307 ---
\setcounter{footnote}{0}%
\opage{307}forget what spirit it is that is in question.
It is not a half-conscious, obscure human spirit that
here stands powerless, groping, before external nature
and confuses its dream images with real facts;
no, it is, by the presupposition of faith,
the almighty Spirit of God, the Spirit of life and of truth,
who is Lord over all reality, on earth as well as
in heaven. When this Spirit in its Word brings the
real under the figurative, the
believer holds himself assured that this does not happen from ignorance
or obscurity, but in the name of truth, for
man's sake, because man, who needs
knowledge of God long before he has come so far on the road of civilization
that he can scientifically clear up for himself his
knowledge of nature, is able to grasp the divine thought deeply and
purely only in such a figurative language. Herewith it is also
presupposed that the spirit, insofar as it
serves the revelation of truth, can also
use the real in another way, that in material
reality it can set a reality of
higher order, that it can, in other words, reveal
itself by signs and miracles.

The question of the signs and miracles of revelation
has --- thanks to the theologians' scientific collusion
with plain atheism --- in recent times become so
distorted that the cultivated person is very loath to enter upon
further discussion of this topic, which is, in his opinion,
as unfruitful as it is worn out. For the question
contains two absolutely dissimilar questions,
the one about God's almighty will, the other about
the immutability of the natural laws. By putting the latter
% --- p. 308 ---
\setcounter{footnote}{0}%
\opage{308}question first, appealing to natural science,
but keeping silent about its well-known concession that our knowledge has an
absolute limit, it is an easy matter to transform
the great fundamental question, decisive for life and death, into
a toothless, theological old wives' question, which nobody
any longer cares to hear.

The considerations and concessions which in this respect
are owed to natural science are, so far as we are able to
penetrate into the matter, carefully set forth in the Philosophy of Religion\footnote[1]{Cf.\ the Author's \emph{Philosophy of Religion}, pp.\ 150---174; pp.\ 476---477, and other passages.}, and we shall not dwell longer here on
this point. What we in this connection must
lay stress on is that the question of signs and miracles
must unconditionally be referred to the question of the
divine will. Is there really an almighty,
creating God, and does this almighty, creating God
have a will? On the answer to this simple question, with
a simple Yes or No, the whole inquiry depends.

That in philosophy there are and must be a multitude of
different shadings in the conception and determination
of the concept of God, we shall be the last to deny;
but from the standpoint of reality, on the stage of life, of existence
and of moral consciousness, it stands immovably
fast that whoever denies God a will denies
God, for a God without will is no God.

The whole account between the faith in revelation
and philosophical atheism is easily settled.
``Will you, the modern Ego of the nineteenth century, shove the will
out of your concept of God and live out your life from
% --- p. 309 ---
\setcounter{footnote}{0}%
\opage{309}mankind's metaphysical principle, very well! that is your own
affair. Surely no one is going to make an attempt
either to force you to faith or in any respect
to thrust upon you biblical representations of the supernatural.
But then you ought also to have so clear a consciousness
of your standpoint that you do not stand there
imagining, to yourself and to the world, that it is a scientific
criticism of the principle of faith that has grounded your
denial of God, whereas, conversely, it is your denial of God
that forms the foundation of your dissolving criticism
of positive religion.'' With this brief retort a
skilled teacher of religion can dispatch the atheist without further trouble.

As regards the individual miracles in the biblical history,
all who agree in acknowledging the
first article of faith must also agree in acknowledging
the possibility of these miracles; for the unconditional presupposition of faith
is that ``with God nothing shall be impossible.''
But from the fact that all miracles are assumed to be possible it
does not yet follow that believers are bound
to ascribe to them all, without distinction, the same outward
reality. It is quite possible that a sign, a
miraculous deed, which according to our way of representing things has the appearance
of having formed itself for visualization in an inward
manner, has under the influence of the spirit been
conceived and presented as if it belonged to the series of
outward events. The main point is that Spiritual Reality
embraces the outward just as much as the
inward, and that every reasoning separation of the
material and the ideal falls outside the interest of faith.
This can be elucidated by a single example.

% --- p. 310 ---
\setcounter{footnote}{0}%
\opage{310}Did Christ really walk on the sea? To such a
question, in so far as weight is laid exclusively on
the outward historical reality, the answer must be short and definite:
there is no human being on
earth who knows; for of such a thing there is no historical
knowledge. World-historical events, which all
proceed from natural causes, science can confirm
by unimpeachable testimonies; but about supernatural
events no scientifically decisive
certainty can possibly be given, for these events
must, by the nature of the case, look highly different,
according as the eyewitnesses have seen them in the light of faith or of
unbelief.

Christ's walking on the sea is no reality for
science, but it is a Spiritual Reality, a reality
for faith. ``For,'' says the believer, ``how
such an event is to be brought into agreement
with the natural laws, of that I know nothing, of that
nothing can be known; but this I acknowledge in faith, that
for God all things are possible, and with that all
reasoning falls away. What appeals to me in this picture of life,
what exclusively captivates my attention, is
the spiritually exalted in Christ's being, his divine
elevation above the outward barriers of nature, and then
first and last the helping hand with which he seizes
the sinking Peter.'' --- This religious point of view is
in principle absolutely different from the theological.

The theological interpreter is a man of science; he
thinks to himself: ``if only I could get sufficient certainty
that it really happened as is told; if I could
by the help of a little criticism, hermeneutics and exegetics
% --- p. 311 ---
\setcounter{footnote}{0}%
\opage{311}get behind the spirit and lay hold of the naked reality,
then I would know what I ought to believe.'' Will
the theologian succeed? Quite excellently! In order not to be
dazzled by Spiritual Reality, he really goes both
behind the spirit and around the spirit; and what does he find
then? A legend, a myth, a distorted reality,
that is to say, a reality in which he cannot believe, so
long as he stands at the lectern; for once he
gets into the pulpit, why then the congregation
will surely be told that he is the man to believe in the naked
--- and yet undistorted --- reality.

What holds in general of the doctrines of faith
holds also of the historical facts in which
they are made visible; they are revelations of the spirit;
but just for that reason these facts speak
not to the sensuous in man, but to
man's spirit. Sensuous reality is reality
for the senses, the reality of what belongs to the understanding is
reality for the understanding; but Spiritual Reality is
reality for the spirit. In the sacred history the Word
is a divine fact, and the fact a
divine Word; the theological mania consists precisely
in wanting --- for the sake of scientific rigor --- to separate
what God has joined together. By such a separation
opinions must necessarily gain the upper hand over
faith, and uncertain conjectures about the outwardly real
furnish the standard by which the spirit is to let
itself be judged. But of this Paul's words hold: ``the sensuous
man does not receive the things that belong to the Spirit of God,
for they are foolishness to him, and he cannot know
them, for they are spiritually discerned'' (1 Cor. 2:14).

% --- p. 312 ---
\setcounter{footnote}{0}%
\opage{312}And now to our next main point:

2) \emph{How can it accord with the infallibility
and unconditional truthfulness of the Spirit
to present events in the sacred history
otherwise than as they actually} happened?

The objections which can be drawn from this against
the concept set up by us, Spiritual Reality, have at
first glance something particularly striking.

``When,'' it is said, ``God in his Word really
wanted to reveal to men how things originally
went with the coming-into-being of the world and the beginning of human life, then it is
after all quite incredible that the Spirit should have presented the matter otherwise than
it in reality stands, that it should have intentionally
--- for from ignorance it could not be --- have
set down something which is notoriously incorrect, that, in
other words, the Spirit of truth should not have told the
full truth.''

``It is, after all, a strange self-contradiction to have to
let oneself be instructed by the Spirit, in so far as it speaks of
the miracle of creation, but in the same breath to be
wiser than the Spirit, in so far as it speaks of the creation
in six days; to have to bow before the truth in the
revealed Word, in so far as it imparts to us a knowledge
otherwise unattainable of the original guilt of the race,
but yet out of respect for geology, botany and zoology
to reserve for oneself the self-satisfied opinion that in reality
the history of the Garden of
Eden, of the state of innocence as a real prehistoric
condition, and the like, is not right. No, honesty lasts
% --- p. 313 ---
\setcounter{footnote}{0}%
\opage{313}longest.'' When one truly wants to be honest and to speak plainly,
then here there is only a choice between One of Two: either to acknowledge
the infallibility of the Spirit in One and All down to letter and
tittle or to reject revelation and build on reason;
for a spirit that is infallible in the one thing, but
not in the other, a spirit that step by step limps
between the true and the untrue, such a limping,
half-erring spirit cannot at any point lay claim
to unconditional faith.''

There is an old, ingrown self-deception lying
at the bottom of these pretended considerations of honesty;
it is this self-deception, not honesty, that we want to get at.

It must stand immovably fast for every honestly believing person
that the Spirit in its Word is absolutely infallible; it is
precisely in order to awaken the recognition of this that we so emphatically
declare ourselves against every reasoning separation
between the thought and the clothing. But the clothing
is, in the name of the Spirit and of truth, equally infallible,
whether the garments, if we may put it so, are
drawn from the lower or the higher, the sensuous
or the supersensuous, the earthly or the heavenly
reality, whether the clothing is woven of immaterial
or of material threads.

The pious self-deception aimed at here consists
in this, that man makes a main thing of what the Spirit in
its wisdom has set down as a side matter, that the believer, without
himself knowing it, places the decisive element in something in which
it is not to be placed, and self-willedly prescribes to the Spirit
an infallibility which it will not acknowledge as its own. Those
who have some acquaintance with the history of the Church and of theology
% --- p. 314 ---
\setcounter{footnote}{0}%
\opage{314}have examples enough of errors and blunders
in this direction\footnote[1]{Cf.\ \emph{On ``the Good Will''} as a \emph{Power in Science}. Copenhagen 1867. pp.\ 40---75.}.

To secure the Old Testament as a perfectly
reliable document of revelation, Protestant
theologians in the seventeenth century proceeded from the
presupposition that even the vowel points were due to inspiration.
They held fast to the assertion that an infallible
revelation required an infallible text; but how uncertain
the text is, when one has only consonants and can arbitrarily
add the vowels, is obvious of itself\footnote[2]{Thus from, for example, the passage in Genesis (ch.\ 23:4), where the meeting between Esau and Jacob is mentioned, one will get out, according as one changes the vowel signs, both that Esau kissed Jacob and that he bit him.}. The
learned Buxtorfs of Basel, father and son, fought valiantly
for the infallibility of the Spirit, in that they fought for the antiquity
of the vowels; but linguistic-historical facts do not let themselves
be reasoned away. The outcome of the dispute showed that the orthodox
zealots did not have the Spirit on their side, that they
on the disputed point had formed an arbitrary
and false concept of infallibility.

The theological dispute about the purity and classicality that has
been wished to be ascribed to the Greek language in the New Testament
stands in a similar way. ``It is the Spirit who in a supernatural manner
has inspired the Apostles with the language; the Greek of the New Testament
must therefore be classical Greek. If there is anyone
who, in spite of Plato, Herodotus and Xenophon, knows the Greek
language from the ground up, then it is surely the Holy Spirit.''
% --- p. 315 ---
\setcounter{footnote}{0}%
\opage{315}To call this assertion in question was, in the opinion of the purists,
to call the infallibility of the Spirit in question. When
it has nevertheless now been established by sober, thorough, many-sided linguistic research
that the Greek of the New
Testament is not classical, but Hellenistic with noticeable
admixtures of Hebraizing idioms, what
does that prove? It proves that the purists have missed
the point of view for the right conception of the gift of language and
sought perfection where it is not to be sought.

The same self-deception repeats itself in the judging
of the historical accounts. As late as the latter
half of the eighteenth century the theologian
Chr. W. F. Walch\footnote[1]{Krit. Nachricht v. d. Quellen der Kirchengeschichte. Leipz. 1770. p.\ 70.} expresses himself on the infallibility
of the sacred writers in the following way: ``they have not only
not erred, but they have'' --- namely on account of inspiration
--- ``not even been able to err, and their narratives
are therefore not merely as certain as any
other historical truth, but even unerringly
certain.'' If the historian here were thinking solely
of the Word as word of revelation, i.e.\ of the truth-content
in its indissoluble unity with form and
clothing, thus precisely of what we call Spiritual Reality,
reality for the spirit, then scientific
criticism, however much it might have contributed
to altering opinions and conjectures about the outward
and finite in the composition of the narratives,
would not have been in a position to alter a tittle in the Word
itself\footnote[2]{On revelation and sacred history. Cf.\ Philosophy of Religion, pp.\ 24---28.}. But that is not how it stands. Instead of
% --- p. 316 ---
\setcounter{footnote}{0}%
\opage{316}applying the measure of the Spirit to the real, he proceeds
backwards and applies the measure of reality, of the
outward finite reality, to the Spirit.
He insists on the infallibility which he himself foists upon
the Spirit, and has no inkling that precisely thereby he gives up the Word
to dissolving criticism\footnote[1]{A couple of examples. In the Second Book of Samuel (XXIV) the following is related: ``And the anger of the Lord was kindled again against Israel, and he moved David against them, saying: Go, number Israel and Judah!'' The parallel passage to Second Samuel (XXIV) is, as is well known, First Chronicles (XXI), which begins thus: ``Satan stood up against Israel and moved David to number Israel.'' Taken as narrative, the one of these accounts stands in literal contradiction to the other. But in the word of revelation itself there is no contradiction; for the moving, which according to the first account proceeds from the anger of the Lord, is referred in the second account to the evil will which precisely in the relation of wrath places itself between Jehovah and David. Thus in John 13:26---27 it is surely both Christ and Satan who move Judas to the betrayal; but who will see in this a religious-psychological contradiction.}.

If we now turn our attention again to what
has been adduced above about creation, Paradise
and the Fall: then all the misunderstandings belonging to the self-deception
will gather in full illumination. The simple
believer, who takes everything according to the letter, because he does not
form any representation that any other conception
was conceivable, is not to blame for the misunderstandings.
The misunderstandings, this must expressly be maintained, begin
only at the turning point where the reasoning
understanding seeks a foothold in the real and, by
setting the natural outwardly together against the supernatural, wants to
control the Spirit.

God's Spirit is then set outwardly against man's
spirit; God is represented as speaking to man in the same
% --- p. 317 ---
\setcounter{footnote}{0}%
\opage{317}sensuous manner as when one man speaks to
another. The Spirit dictates to the holy men what they
are to write; if there are errors in the dictation, then the Spirit
must of course bear the blame: it is the strongest caricature
of Spiritual Reality that can well be conceived.
The fundamental error is that the supernatural is treated
as if it were natural.

In revelation the relation between God and man's
spirit is a relation of spirit; the Spirit speaks in an inward,
not in an outward manner; even when the impression
is so strong that the Word sounds to the ear of faith, the outward
nevertheless springs from the inward. Therefore the Spirit speaks
a human language; it expounds its Word precisely in
the language which the people understand to whom it
speaks; therefore it spoke Hebrew to the Jews, Hellenistic
Greek to the Greeks of that time, and would to this very day
speak Danish to Danes in the first instance, German to Germans
and French to the French.

If the relation between God's and man's
spirit is conceived as a real relation of spirit, then one will also
be able to understand that the Word of revelation does not merely clothe
itself in human language, but in general adapts
itself to man's circle of representation, makes use of
his own view of the world, his natural view of
heaven and earth, in order by these means of connection to
make vivid the truth on which alone it depends, % print: kvorpaa (as printed)
the truth of faith, for spiritual awakening, for peace and
salvation.

From this point of view, should it not now be
obvious what self-deception lies in the demand
that the Spirit in its Word about creation, Paradise and the
% --- p. 318 ---
\setcounter{footnote}{0}%
\opage{318}Fall should have laid at the foundation the astronomical, geological and
natural-historical representations which prove to
hold good in the more recent science! We ought indeed
to be ashamed of so willful, so unreasonable a decision
as this: either it is the Spirit that has erred, and
then the Word is not credible, or the Spirit itself must
be infallible also in descriptions of the physical, and
then we must see to it that we instill in dear youth every
possible distrust of science.

It is written: be simple as doves;
but it is nowhere written: be stupid as
flounders. Everyone who in our time will use his sound
understanding and not, in spiritual arrogance and foolish defiance,
cast aside human knowledge, will, as regards the structure
of the universe, easily convince himself that the offense
at the Copernican system rests on a
self-conceited imagination. Let it be a thousand times true
in the astronomical sense that it is the earth that goes
around the sun, it is then just as true that the Spirit of revelation
cannot possibly have begun by hinting
at this, without making its Word wholly unintelligible to
men.

Let it be well noted that enlightenment about so
complicated a matter is not given by a single remark.
The mere having to shift one's standpoint from the earth
to the sun and to determine the heliocentric places of the stars
--- what an exertion it costs the beginner. There
are among the so-called Cultivated thousands upon thousands
who ``believe in'' the Copernican system, but who would
fall into the most painful embarrassment if they were to give
an account of the proofs. And then the Spirit was supposed, moreover,
% --- p. 319 ---
\setcounter{footnote}{0}%
\opage{319}to have proved its astronomical infallibility by making
intelligible to Abraham and Moses what has scarcely yet
become intelligible to the thousandth part of those who
drift with the stream of cultivation.

In its Word the Spirit is infallible; for the revelation
of truth in and through the Word it unconditionally takes all
reality, both the outward and the inward, into
its service; therefore the eye of faith cannot possibly see the real
before it in a merely natural light.

Outside revelation and faith there is no Paradise
and never has been; but in the Word of revelation, for
the gaze of faith, Paradise has full reality. The believer
sees everything in the holy light, and everything is real: the Garden
of Eden, the Tree of Life and the river Gihon; he sees both
the serpent and the woman. But in the light of science there is
nothing here to see, nothing at all.

From a geographical vantage point one beholds the Desert of Sinai:
granite mountains with wild, wondrously jagged peaks rise
above water-rich, fruitful valleys; one sees Horeb,
Jebel Musa, the true peak of Sinai; it is all
real: here is reality enough for geography.
But to this kind of reality faith is blind; in
the eyes of the geographer it is like the owl, which cannot see by
day. This comes from the fact that the eye of faith is differently
organized and needs a different daylight from the observer of nature's.
Only when the light of the Spirit shines upon Mount
Sinai does faith become seeing; then it really sees the
top of the mountain; it sees Jehovah and Moses: in the thundercloud is
Jehovah's finger and the tablets on which he writes!
But, curiously enough, now the student of nature can see nothing;
now it is he who is blind as an owl.

% --- p. 320 ---
\setcounter{footnote}{0}%
\opage{320}That the Sea of Gennesaret is a blazing mirror when the sun
rises and the rays glow on Mount Tabor, is a
landscape glory which has never yet been seen with the eye
of faith. This by no means signifies that the believer should be
shy of light or lack a sense for the beauties of nature; but it
signifies that faith as faith must have its own light.
When this light is kindled, it sees reality everywhere,
the natural as well as the supernatural; it sees
the mountain and the Transfigured One on the mountain; it sees in
the dusk of night Gennesaret by the sea, and the waves and
the ship and the Savior walking on the sea. It is Spiritual Reality,
reality for the spirit.
\signature{R. \emph{Nielsen}.}

\clearpage

% the title page of Bind 5 (1871-2-for-idee-og-virkelighed.pdf), in English
\begin{titlepage}
\begin{center}
\vspace*{3em}
{\Huge For Idea and Reality.}\\[1em]
\rule{8em}{0.4pt}\\[1em]
{\large A Journal,}\\[0.8em]
{\small edited}\\[0.4em]
{\footnotesize by}\\[0.6em]
{\large R. Nielsen,}\\[0.4em]
{\large Bj. Bjørnson and Rud. Schmidt.}\\[1em]
\rule{8em}{0.4pt}\\[0.6em]
{\large Year 1871. Second Volume.}\\[0.6em]
\rule{8em}{0.4pt}\\[4em]
Copenhagen.\\[0.4em]
Otto B. Wroblewsky, Publisher.\\[1.5em]
{\small Printed by G. S. Wibe.}\\[3em]
{\small Translated from the Danish}
\end{center}
\end{titlepage}


\essayhead{More on Spiritual Reality,}{occasioned by the Contributions of Pastor V. Birkedal and Pastor C.\\ Møller.}
% --- p. 1 ---
\setcounter{footnote}{0}%
\opage{1}Anyone who, with a collected mind and an unprejudiced eye, sets out to go through
the accounts in the Old and the New Testament
and to make the supernatural events
present to himself must surely, unless he lacks all
the presuppositions of cultivation and the conditions of independent thinking,
come to have strange thoughts. How different the whole
real order of things looked in the days of the prophets,
of Christ and of the apostles from the way it looks in ours!
It is as though Holy Scripture spoke of another
heaven and another earth, another world and another
God, in short: of an altogether different reality, not merely inward but
also outward, from the one we actually know.
The present age, which spies and tests and is otherwise so sure in its observations of reality,
even if it strains the powers of its reason
to the utmost, is not able to bring
its concept of reality and of real facts
into harmony with what the Bible calls the real. To
take refuge in science for help and assistance
certainly has its significance, when the task is to be rid of the
% --- p. 2 ---
\setcounter{footnote}{0}%
\opage{2}biblical reality; for how representations of the
supernatural can be dissolved into legends and myths, of that
science has excellent knowledge, and no one to whom
the miraculous events are a burden shall regret
having sought its guidance. But to want --- in
the conceit that faith is of one kind with knowledge --- to plague
science into standing surety for the dependability
of that miracle-business which is so offensive to
``all the cultivated,'' can of course occur only to those who
still have a theological screw loose.

If the natural events in profane history
and the supernatural ones in the sacred history are really
to be seen in the same light and subjected to the same criticism, then
Christianity is lost to civilization: the credibility of the Bible
cannot be saved. To get out of
the confusion in which an objective, merciless criticism in the name of science
has entangled the unhappy theology, so that
it knows neither out nor in, dares to be neither critical
nor uncritical, neither give up the miracle nor hold on to
the miracle, there is only one way out: unconditionally to acknowledge
the independence of the religious principle and to declare
the sacred history to be of a different kind, \emph{absolutely} of a different kind, from
the profane.

Once it is fully seen that the biblical accounts
must, by the nature of the case, be of another origin,
be grounded on another reality and fall
under another field of view than accounts that belong
in a general world history, it follows
indisputably from this that the biblical accounts must be seen
in their own light and judged from their own standpoint.
As reality is, criticism must be. In order then to
% --- p. 3 ---
\setcounter{footnote}{0}%
\opage{3}give that higher reality a designating name I have
called it Spiritual Reality, reality for the spirit, in that
by the first expression I have aimed in particular at the divine origin of the
facts, by the latter, on the other hand,
at the consciousness that is to make sure of the
truth of the facts. I shall now seek to bring
the main thought thereby indicated as close to the eye as
possible.

\begin{center}I\end{center}

\emph{Spiritual Reality supplies the first necessary
presupposition for the judgment of the
sacred facts, when the religious
consciousness has reached the stage of cultivation} at which
\emph{the rational objections of science can}
make an impression upon \emph{it, so that it must not merely outwardly},
\emph{but also inwardly find itself in conflict with scientific criticism}.

``Reality for the spirit,'' says Pastor Birkedal, ``is
not the first reality.'' Reality must ``first
prove itself to another capacity of reception (organ) before it
can prove itself to the spirit, if Christianity really
is to become our salvation, our life. Certainly it must
happen through the spirit, because Christianity is a realm of spirit,
in which nothing comes and is bestowed \emph{without} the spirit, and
to that extent one may call the reality of Christianity
`Spiritual Reality.' But the explanation of this
expression: `Reality for the spirit' nevertheless seems to me
not to correspond to the `real' relation, when one is to
begin with it.'' Yes, when one is to begin with it!
In this Pastor B., in my view, is perfectly right.
% --- p. 4 ---
\setcounter{footnote}{0}%
\opage{4}But the point is that one has neither begun nor
is to begin with it in that way. As to the origin of the life of faith we
are, as Pastor C. Møller also points out, in
all essentials surely quite agreed. Here, in this
context, the question is not how faith comes
into a man's heart, or how a man
comes to faith; but the question is how
a man who believes from the heart in God's revealed
Word and with this faith has lived himself into the whole
circle of representations of the Bible and found comfort and joy in the assurance
that the sacred, supernatural facts are not
floating legends and images, but fully dependable
reality, how, I say, such a man will
be able to hold fast to this reality in the face of
a criticism which, with weighty scientific grounds,
impresses and in many turns repeats that proposition so well known from
negative knowledge: \emph{The divine
cannot have happened thus; what
has happened thus cannot be divine}.

It is not respect for the facts, not
love for the Word of the Bible or a thorough study of
Holy Scripture, that makes scriptural theology spiritless; the
spiritless, that is to say: the theological, lies precisely in the
unhappy misunderstanding with which one begins. Without
making clear to oneself that it is a quite peculiar kind of
reality one has before one here, one proceeds tacitly
from the unclear presupposition that reality is reality,
event event; --- a materialist is after all a materialist,
as the vaudeville has it. Scriptural theology
sets, in other words, the reality of the supernatural,
mark well: as regards the criterion of reality,
% --- p. 5 ---
\setcounter{footnote}{0}%
\opage{5}on a par with the reality of the natural, and
what is the consequence? The same historical proofs that
are applied to defend the natural events are applied
now also to defend the supernatural ones; the same
calculation of probability that is brought to bear on the former is now
also supposed to fit the latter. Then the biblical theologian
appears with what he calls ``a believing mind''; that is to say: he
is willing to concede that something wonderful has happened,
but wherein the wonderful really consists is to be
found out only by continued scientific investigation; he
is willing to accept miracles, provided they are not too
absurd and improbable to him. But unfortunately for him, there
is in the whole Bible not a single miracle, no, not a
single one, which is not, from a scientific standpoint, closely and
narrowly regarded, absurd, improbable and repellent.
Scriptural-theological criticism assumes a devout mien
when it is to defend the great facts of the resurrection and the ascension,
whereas it wholly keeps its
scientific bearing toward the miracle of the Gadarene
swine; it thinks as follows: with swine and devils
one need not pay compliments. But where is
faith now, and where is criticism?

That the miracle of the resurrection has a quite different, pervasive
significance for the religious life than the event
with the possessed man at Gadara, only a madman
could think of denying; but here the question is not
about a miracle's greater or lesser religious significance;
the question is about the reality of the miracle and the greater or lesser
probability of miraculous
reality. When this point of view is held fast, and it must be held fast strictly,
if criticism is to mean anything at all, then
% --- p. 6 ---
\setcounter{footnote}{0}%
\opage{6}not much acuteness is needed to see that
all miracles are equally probable, that is, equally improbable;
they resemble one another unconditionally in One thing: they are all miracles.
The miracle of the resurrection does not come a hair's breadth nearer
to ordinary historical probability
than the miracle of the demons in the herd of swine.

The more the biblical theologian wants to light up, with the torch of knowledge,
the dim ground that bears the sacred facts,
the darker, the denser the mists that rise from
the depths, and the more uncanny he must feel.
There is in this surrounding everywhere a Something that goes against his
consciousness, an Inexplicable that presses upon his
chest and weighs on his scientific conscience
--- if indeed he really has a scientific conscience.
Hence the boundless irresolution in the face of
what is most decisive; hence the endless wavering between acceptance
and non-acceptance, the benumbing halving and
meaningless explanation that something has certainly happened,
something extraordinary; but what it is that has happened, that
one does not know. So the region of revelation becomes a land
of mists, its content a billowing element in which
one goes to the bottom in the diving bell of science.

Theology must go. Here another starting point must be chosen;
here one must begin from another ground.
It must be settled once and for all that he who steps
into the light-region of revelation and feels himself addressed by
the supernatural phenomena has crossed the boundary
of knowledge and said farewell to the compelling, universally necessary
grounds of reason. In so doing the believer has by no means
either given up the laws of thought or extinguished the light of reason;
but One thing he has given up, the false justification, the
% --- p. 7 ---
\setcounter{footnote}{0}%
\opage{7}justification in objective knowledge. He feels and understands that
religion could not possibly be religion if it did not
have its own origin, its own self-valid principle;
he feels and understands that facts which proceed
from a supernatural principle (the holy, almighty will)
must bear a corresponding supernatural stamp of reality;
he feels assured that the whole self-contained
circle of cognition with which his soul is familiar, and
in which his better being feels at home, belongs to a
reality that can no more be confirmed than undermined
by science. With this foothold in the
supernatural distinguishing marks of reality he is, as a believer
and yet as one who knows, prepared to receive the
objections of scientific criticism. These objections
will be no more able to shake his faith in
revelation than to weaken his respect for science.
He is so far from becoming embittered in the name of religion
at the attacks of objective criticism that he
will much rather see that, with its presuppositions,
its foundation and its measure of reality, it cannot possibly at any
point appropriate the supernatural, and that in this
respect it is legitimately excused. Gradually, on the whole as
in particular, it will become quite evident to him
how the whole web of critical difficulties is essentially
due to theological weavers, who have made the beginning
by weaving together two kinds of reality, physical
reality and Spiritual Reality.

\begin{center}II\end{center}

\emph{Spiritual Reality expresses the concept of reality
in its whole extent and
% --- p. 8 ---
\setcounter{footnote}{0}%
\opage{8}thus embraces both the inward and the outward
real, both bodily and incorporeal
reality, but always with particular emphasis
on the spirit itself, the almighty will},
\emph{as cause. From the standpoint of religion and faith
all reality, the natural just as
well as the supernatural, is in the last and highest
instance to be brought under Spiritual Reality},
\emph{reality for the spirit}.

If the sole unfailing mark of reality
were the outwardly sensible, the properly tangible,
then earths and stones would presumably have to be
placed in the highest position in the realm of reality. We should then get
a materialism of the very coarsest kind; for the materialists
are otherwise wont to let everything real rest on a
swarm of minute parts (atoms) that are not directly
tangible. Reason and religion agree in this,
that there is working not only in the corporeal but also
in the incorporeal, and that as a consequence there
must be an incorporeal reality just as well as a corporeal one.
Those who concede that there is reason in things must then
also concede that everything corporeal has its ground in something
incorporeal, unless one wants to regard eternal
reason itself as a body, a kind of invisible matter.

The conflict between religion and reason first comes to
light when the invisible cause of reality is to be more closely
determined. Physics keeps to the forces of matter
and researches laws; it has to do only with the natural
side of reality, thus only with natural reality, not with
Spiritual Reality. Philosophy strives to
present real nature, the in itself
% --- p. 9 ---
\setcounter{footnote}{0}%
\opage{9}coherent natural reality, as grounded in the essence of the spirit;
but in such a way that the working of the spirit in nature at no
point exceeds the boundary of the conditions of nature.
Religion --- we speak here of revealed religion ---
on the other hand impresses upon us, as its very first principle,
that God's free, almighty will is the original cause
of all reality. \emph{In the beginning God created
the heaven and the earth}; such, as is well known, are the first
words of Holy Scripture; by this it is expressly said
that the reality of the whole world is at bottom Spiritual Reality.
To believe in the first article, on creation, and
yet to deny that Spiritual Reality is the true, comprehensive
expression for all reality is a contradiction.

The relation between nature and creation can hardly be characterized
more clearly than by saying: creation in its
ground and coherence is \emph{nature}; nature in its origin
from and dependence on God is creation.\footnote[1]{What matters here, as everywhere, is to hold the points of view sharply apart from one another. To explain the natural half-and-half as supernatural and the supernatural half-and-half as natural is nonsensical; by mixing together physics and religion and mediating the mixture in dogmatic speculation one does religion and science a poor service.} The confirmation
that created things have their nature is
to be seen in the laws; the sign that real nature was originally
created is to be seen in the miracle. The miracle is simply
to be defined as the sign of omnipotence, an express
declaration that God really has a free will, and
that this will has all things unconditionally in its power.
If real nature were not receptive to the miracle,
then the realm of nature would not have been created: to deny the possibility of miracle
% --- p. 10 ---
\setcounter{footnote}{0}%
\opage{10}is to deny the possibility of free creation. The
real miracle, whether it otherwise takes place in the
inward or in the outward, in the incorporeal or in
the corporeal, is then in the very strictest sense to be called
Spiritual Reality, reality for the spirit.

Holy Scripture everywhere places, as regards reality,
the miracles in the spiritual on a par with
the miracles in the bodily; this is a main feature that
must by no means be overlooked. Some examples. Jehovah
reveals himself to Abraham, first in a vision (1. Mos. ch. 15),
next (ch. 18) in an outwardly real form;
the first revelation thus points to the incorporeal,
the last to the corporeal. Who will now, from
the standpoint of faith, maintain that the first manner of revelation
is less supernatural, less real and therefore
less dependable than the last? That the last goes
deeply into the outwardly natural is seen from the hospitality
with which the patriarch entertains the three heavenly ones:
``he took thick milk and sweet milk and the calf that he
had prepared, and set it before them, and himself stood
before them under the tree, and they did eat.'' Is it now permissible
to say that the first reality, the reality in the vision,
is only a trifle of Spiritual Reality, a fleeting idea-reality,
perhaps a waking dream, an imagining of Father
Abraham, but that the last is the proper, true and
solid reality, because it is a reality with calf's flesh
and thick milk, and thus a reality that falls outside
the spirit?

With regard to the concept of reality itself there is
no difference between the Old and the New Testament.
We learn it from the miracle of the resurrection. That Christ really
% --- p. 11 ---
\setcounter{footnote}{0}%
\opage{11}has risen from the dead, that his revelation to his own
has taken place not merely \emph{inwardly}, as in a vision, but also
\emph{outwardly}, indeed (John 20:27; cf. Mtth. Mr. Luc.) in
a quite \emph{tangible} way, is as strongly attested as
it well can be by believing witnesses; but no one will surely
deny that the body with which Christ was crucified, and with which he rose, has undergone through the resurrection
such a transformation that he is now recognizable,
now again unrecognizable, that he comes as a vision and
yet as one who has flesh and bones and an appetite for food
--- they gave him a piece of a broiled fish and of
a honeycomb (Luke 24:42) --- that, in other words,
his bodily reality has taken up the natural into
the supernatural, that it corresponds in one and all to Spiritual Reality,
reality for the spirit. The theological men, who
cannot be satisfied with the spiritually real, but
must of necessity have either ``honeycomb'' and ``broiled
fish'' or ``calf's flesh'' and ``thick milk'' thrown in, to
make good the spiritual and fill out reality, will
with all their theology have trouble enough in believing
in Christ's real resurrection.

\begin{center}III.\end{center}

\emph{Spiritual Reality is the categorical expression}
of \emph{the unity of the form of revelation with the
revealed content. It is the at once
supernatural and natural unity of spirit and
word that makes the word God's Word and the Word of God
itself the highest reality, the real} facts' bearer and presenter. The
\emph{presenting word does not have the reality in question
% --- p. 12 ---
\setcounter{footnote}{0}%
\opage{12}outside itself, but in itself.} If
\emph{the word is separated from its spiritually real} content, then \emph{the
linguistic, the grammatical becomes the main thing:
the revealed word is transformed into a
philological text, and strife arises between}
spirit and \emph{letter}.

``In the beginning was the Word, and the Word was with God,
and the Word was God'' (John 1:1). --- Were the task
to state briefly and precisely the point of view of the word-picking,
dyed-in-the-wool philological exegesis, which turns the Word of God into
sentences and the doctrine of sentences into grammar, it could
best be done by a suitable paraphrase of the passage cited:
``In the beginning was Grammar, and Grammar
was with God, and Grammar was God,'' and so on.
From an objectively learned standpoint it will hardly occur to
anyone to cast aspersions on a deep-going philological
study of the original text of the O. and N. Test.; nor
shall it be overlooked that the yield of such a study
mediately affords guidance also to those who
solely in a religious interest wish to secure an understanding of
the biblical text. But in the name of religion a protest
must be lodged against the minutiose-philistine importance
and ever-increasing prolixity with which biblical philology
in recent times has contrived to make its exegetical analyses
into a scientific necessity of faith, a kind of
grammatical concern of salvation. So it happens that
the side issue becomes the main issue, and the main issue the side issue; so it happens
that the Word of God is dissolved into sentences and letters.

The relation between the Word of God and the biblical text
may be briefly elucidated in the following way. In the gospel of the resurrection
the believer possesses a perfectly true and
% --- p. 13 ---
\setcounter{footnote}{0}%
\opage{13}dependable presentation of a real historical event.
The miracle of the resurrection is God's work; the infallible
account of this miracle is itself a miracle, is
likewise a work of God: the event and its apprehension,
the outward and the inward reality, fuse in the Word of God
into one single indivisible reality, Spiritual Reality.
And yet there is a difference between the description
and that which is described. The evangelical account is a
mirror that gives back the individual features of reality in a faithful
image. Faith's eye looks in the mirror at the image and in
the image at the reality, the real event:
the gaze is to that extent turned outward. If now, in the outward
sense, there were only one gospel, one account,
then the outward-turned gaze would involuntarily become inward-turned,
the unity and harmony between the fact and the apprehension,
between the outward and the inward real,
would then remain unaltered. But the divine
cause has, and by virtue of the nature of revelation must of necessity
have, its human side. The one gospel
is divided into four gospels, the one infallible account into
four accounts differing from one another and yet infallible.
If one errs about this difference, there arises
an outward separation between the Word in divine
truth and the account in human form, a worrying
divorce between the Word of God and the biblical
text; there will be something for Grammar to do.

In Matthew (ch. 28) it begins with Mary
Magdalene and the other Mary coming, as it dawned toward
the first day of the week, to see the grave. What
a sight! The angel of the Lord descends from heaven to roll the stone
from the grave; his appearance is like lightning; the watch flees
% --- p. 14 ---
\setcounter{footnote}{0}%
\opage{14}and the angel addresses the women. But none of the three
other Evangelists seems to know anything of a guard
having been set at the grave; nor, according to their account,
was any of the women an eyewitness to how
the stone was rolled away; on the contrary, it had already
happened before they came. Yet another discrepancy. Mark
(ch.\ 16, v.\ 1) says expressly that the women, only when
the Sabbath was past, bought the spices for the embalming;
but Luke (ch.\ 23, v.\ 56) testifies that even before the
Sabbath they had prepared spices and ointments: how could they
have prepared the spices and ointments that they had not yet bought?
Even if one can, by a philological treatment of
the verb ``\emph{prepared}'' (ἡτοίμασαν), glide somehow
past this apparent contradiction, how does
one get through the difficulties of the contradictory
accounts of the angels? Matthew and Mark
bear witness that there was only one angel at the grave,
but Luke and John state expressly that there were two, and yet
their testimonies do not agree outwardly either. According to
Luke the angels are seen by all ``the women from Galilee'';
but according to John no one else gets to see them but
Mary Magdalene. Still more striking is the different
manner in which the angels express themselves. In
Luke they deliver a detailed, instructive discourse to
the women; in John the address to Mary
Magdalene confines itself to the question: Woman, why weepest
thou? According to Matthew and Mark, Christ commands immediately
after his resurrection, both through the angel in the grave
and through his own spoken word to the women, that the disciples
shall meet him in Galilee; but according to Luke he commands
that they shall all remain in Jerusalem until
% --- p. 15 ---
\setcounter{footnote}{0}%
\opage{15}the Holy Spirit was poured out upon them. By going
into particulars in this way the reader will discover one striking
discrepancy after another. Lessing, in his
defence of his ``Unknown''\footnote[1]{\emph{Eine Duplik}. 1778. Sämmtliche Schriften, 4ter Band. Berlin 1825. S. 95 flgd.}, has demonstrated no fewer
than ``ten contradictions.'' Ten contradictions in four short accounts!
Is it not dreadful?

If what is offensive in such contradictions really
had its ground in Scripture itself, then such a writing ought
not to bear the name: the Holy Scripture; it would be more
fitting to call it the obscure, confused, contradictory
Scripture. But the fact is that the offence aroused by the contradictions
has its ground not in Scripture but in the false point of departure
of the interpreters of Scripture. By beginning with a philological
criticism of the text one unawares tricks faith into the
domain of knowledge and comes to apply a standard
that stands in manifest conflict with the essence of revelation.
The fundamental error is that one absolutely insists on treating the Bible
like any other written document of antiquity and regards it
as a settled matter that in the Holy Scripture, just as
much as in profane writings, a final critical
separation can and must be made between the facts
presented in the Word and the outward supposed
reality of these facts. But in the Holy Scripture it is
precisely a self-contradiction to make this separation. The miracle
and faith are inseparable; take the miracle away, and there
is no faith; take faith away, and there is no miracle.
It is not merely in the later writing-down that it
essentially depends on the faith-witness's peculiar
% --- p. 16 ---
\setcounter{footnote}{0}%
\opage{16}conception and gift of presentation; no, the fact and
its conception are from the very first moment in the germ itself
inseparably united: the outward reality in the event
and the inward reality in the conception is one indivisible
fact, one reality: it is Spiritual Reality,
reality for the spirit.

That the matter really stands so is seen most clearly
from this, that the incidental circumstances that from a critical
philological point of view should speak \emph{against} come
precisely to speak for the credibility of the Gospel. a) The elements of
the outward all belong to one reality, one great, marvellous
event. But the event has several members
and many details; these elementary details are
so dissolved in the appropriation and transformed in the presentation
that what was lived through outwardly and inwardly cannot
be separated in the Word. The indissoluble unity in the Word of the
outward and the inward is everywhere a mark of character
of revealed truth. b) The deviations in conception
which, when one stares at the outward and
clings to the letter, become real, offensive
discrepancies, change into eloquent testimonies
to a living, free, personal appropriation, when they are seen in the spirit
and understood in the Word. Discrepancies in one narrative
with other narratives can undeniably present
hard knots, insoluble contradictions, when the narrator
gives himself out to be copying an outwardly given
reality; for it is literally true that one angel
is not two; that what one buys before the Sabbath one does
not buy after the Sabbath; that Galilee is not Jerusalem, etc.
But in the spirit and the Word it is still more decisively
true that, in relation to an overwhelmingly great and miraculous
% --- p. 17 ---
\setcounter{footnote}{0}%
\opage{17}event, each individual witness must be so strongly
gripped, so wholly filled and inwardly moved by the event
in its entirety, that the human side of the appropriation
will also show itself, so that in reproducing
the main parts and in a subjective painting-out of the
individual features he involuntarily comes to individualizing
alterations. The Evangelist is moved by the spirit
and in the spirit by reality; but he is no photographer
who wants merely to give us a copy of an outward object,
any more than he is a scribe who mechanically writes down
what the spirit mechanically dictates. The Evangelist is
a witness to the truth, of course, not a witness to the truth in the dogmatic-episcopal-modern
sense --- to such a height of
brilliance he can naturally not raise himself ---
but he is nevertheless really a witness to the truth; he testifies
in the spirit to reality: the content of the testimony in the Word
is Spiritual Reality.

\begin{center}IV.\end{center}

\emph{Spiritual Reality is characteristic of a
standpoint from which it becomes possible} to \emph{carry through
the separation between the scientific
and the religious conception of the sacred
history so thoroughly in principle that, both about the
whole and about its individual parts, there can form
two kinds of criticism: criticism of knowledge and criticism of faith}.

Among the many turns of phrase, partly superficial, partly subtle,
with which those who believe in knowledge have sought to
ward off the strict antithesis between religion and science,
belongs also this one, that a religious interpretation of the
Holy Scripture must become groundless, arbitrary and vulgar,
% --- p. 18 ---
\setcounter{footnote}{0}%
\opage{18}when it is torn loose from scientific research. It is, they
say, to cast a shadow upon learning and scholarship
when one makes the common people believe that simple
believers can understand the Gospel with the bare heart
without guidance from genuine scriptural scholarship. To let everything
depend on a living faith and to proceed from the simple
assertion that miracles are miracles, that one must either
accept them all or reject them all, is an altogether too convenient
way of escaping the difficulties; in short;
Theology wants to impress the public by invoking
thoroughness. We must remark here that the swaggering
blasts on the trumpet of thoroughness are only signals of distress
that have nothing to do with Spiritual Reality. We are
so little inclined either to disparage thoroughness
or to flatter self-satisfied thoughtlessness that, on the contrary,
to the best of our meagre ability we are intent on getting at the puffed-up
and on pursuing want of principle to the utmost, whether
it springs from scriptural scholarship or from self-conceit
or perhaps from both together.

It is rightly maintained that there must be a
scientific interpretation of the Holy Scripture alongside
the religious; but it must be repeated again and again
that the two interpretations belong each to its own circle of cognition,
that the one cannot ground the other, but that they may well
exert a mediate, indirect influence upon each other\footnote[1]{It is presupposed, of course, that those readers who wish to acquaint themselves with the way of looking at things expressed here will, as time and occasion allow, make themselves familiar with the writings belonging here: \emph{Obstacles and Conditions for the Spiritual Life}, \emph{The Good Will as a Power in Science}, \emph{Philosophy of Religion}, the essay on \emph{Spiritual Reality} (10th number of the journal, October 1870).}. It is this relation that we shall elucidate by
% --- p. 19 ---
\setcounter{footnote}{0}%
\opage{19}a brief characterization of the criticism corresponding to knowledge and that corresponding to faith.

a) \emph{Criticism of knowledge} is characterized by the splitting
and dividing up of the whole that is locked together in Spiritual Reality.
The fact falls outside the conception; the account
is one thing and the real event another:
agreement is possible, but --- this doubt is
the spur of criticism! --- disagreement is also
possible. By careful comparison of the individual features of every
account before us with one another
and with other corresponding accounts, if such
exist, the interpreter strives to come as near as possible to the unknown
reality lying behind it. The
distinguishing mark of criticism is that the investigator everywhere goes around the spirit
and behind the Word in order to search for the fact,
the outwardly real. What has really happened, then? Was
Christ really tempted by the Devil? Was he really
transfigured on the Mount? Did he really walk
on the sea? Was he really able by his word
to still the storm? etc. Now thoroughness is what counts;
the more thoroughly the philological-critical interpreter proceeds,
the more clearly the yield of criticism shows itself.

The yield is that reality vanishes in a mystical
half-darkness; the natural features remain standing, but the
supernatural are only glimpsed in ``legends and myths''\footnote[1]{Cf.\ An Interpretation of the Sacred History (journal, 4th number, April 1870).}.

\emph{The Temptation in the Wilderness} has no reality, neither
as regards the Devil nor as regards Christ. ``As
% --- p. 20 ---
\setcounter{footnote}{0}%
\opage{20}regards the Devil, the procedure in the whole difficult struggle, the flatness in the ordering and introduction
of the temptations seem as difficult to reconcile with
the shrewdness and calculating cunning that were here supposed to
reach their culmination, as, on Jesus' side, the patience in
continuing these negotiations, and in letting himself be led by the evil spirit
from place to place through the various
stages of the trial, difficult to reconcile with the ideal of
moral purity, which must at once have withdrawn from
contact with the unclean.'' (H. N. Clausen.)

``\emph{The Transfiguration on the Mount}'' \emph{has no true}
reality; here is --- despite the strong illumination ---
only ``a picture of mist in mystical half-light.'' ``The three
eyewitnesses (Peter, James and John) were heavy
with sleep, they knew not what they said, for they were sore afraid;
when they fully came to themselves and looked around,
they saw no man, save Jesus only.'' (H. N. Clausen.)

``\emph{The Walking on the Sea}'' does indeed point to
a real sea and a real walking, but the relation indicated in the Gospels
between sea and walking is
hardly real. Philological criticism has ``grounded
supposition that the episode about Peter is a later addition,
owing to the Greek translation or adaptation
of the Gospel of Matthew, a supposition that is
strengthened by the whole character of this narrative, which
recalls the embellishing and intensifying tendency in
the apocryphal Gospels.'' (H. N. Clausen.)

``\emph{The Stilling of the Storm on the Sea of Galilee}''
has in the evangelical presentation had its
reality so distorted that the fact itself has been
turned backwards. The Evangelists (Matt., Mark, Luke) have
% --- p. 21 ---
\setcounter{footnote}{0}%
\opage{21}namely in their simplicity assumed that it was the storm
that Jesus rebuked, whereas criticism must find that
it is not the storm but the disciples who get a reproof;
for it is far more probable to suppose that Jesus' threatening
words ``were addressed to the disciples rather than to the elements,
to weather and wind.'' (H. N. Clausen.)

It is not merely a certain limited number of individual miraculous deeds,
perhaps less essential to the religious consciousness,
whose reality is in this way
dissolved and turned upside down, for the great facts of salvation,
the fundamental miracles on which the whole revelation
rests, suffer the same fate.

That a criticism that goes around the Word and behind
the spirit cannot acknowledge Christ's conception by the Spirit
is a matter of course. What should move it to
such an acknowledgment? The ``fragmentary
narratives'' of the Prefatory Gospel are of doubtful origin; ``a mythicizing
investiture,'' an ``embellishing setting,'' has
stepped in ``filling out and shaping'' to make up for ``the
scantiness and vacillation'' of the historical traditions.
--- With the end it goes exactly as with the beginning.
The Gospels of Christ's resurrection and ascension
are just as mythicizing, just as fragmentary,
just as embellishing as the Gospels of his birth
and childhood. ``The mutual divergences with which the
four evangelical accounts of the resurrection cross one another'' have so disturbing an effect that ``only
with great uncertainty can a tolerably coherent
narrative be produced by comparing them with one another.''
Still more distressing is the case
with the ascension. That it is attested as a
% --- p. 22 ---
\setcounter{footnote}{0}%
\opage{22}mental fact must be conceded; but in what this fundamental fact
actually consisted, despite the thoroughness of the investigation,
cannot be searched out; to say it in plain words: the Evangelists
know no more, and the interpreters neither.

b) \emph{Criticism of faith} is now recognizable by its holding
the facts fast in the evangelical conception.
It is indisputably shown, as Pastor Birkedal emphasizes,
that the believing heart cannot be content with ``shadow-deeds,''
but needs ``real, tangible facts'';
but how the highly honoured author has hit
upon the thought that I possibly regard Spiritual Reality only
as a kind of Idea-Reality and the revealed facts
as shadow-deeds, he must himself understand,
for I do not understand it. What I, according to the principle of
criticism of faith, must unswervingly maintain is that the believer, with
all his need for ``tangible facts,'' may not
go behind the Word and around the spirit; he must
not imagine that he can look over the Evangelist's
shoulder and, by peeping as out of a window for something
outward, catch sight of a better, surer and more reliable
conception of reality than the one that is given him in
the evangelical Word.

Let us now see how this principal proposition, valid for Spiritual Reality,
can be applied. Pastor B.
gives the matter the following turn. ``The heart says: The Lord's
signs must be real. --- He must really have
been able to command the forces of the world and not need to
let himself be bound by the laws of nature --- that is necessary for
my faith; for if he could not do these miracles in nature,
then I cannot believe in him himself as the greatest
sign, as the one who is Lord of Nature, then I could
% --- p. 23 ---
\setcounter{footnote}{0}%
\opage{23}not believe either that He does the still greater miracle,
in spite of my nature and its laws, which by \emph{their}
necessity must lead to death, of giving me his
life and letting me pass with him over the Sea of Death, so that I
do not go to the bottom in it. I have that by \emph{feeling}:
if he could not bring nature out there to obey,
then neither could he bring my still more rebellious
nature to submit to Him.'' This is, in the
name of faith, both certain and true in every respect. But
if Pastor B. now means that he can have, ``by feeling,''
a sure knowledge that Christ really walked
on the sea, that by his word of omnipotence he really stilled
the storm, etc., then I beg him to consider whether
he has not here allowed himself a hasty conclusion. So long
as he speaks in general terms of what the Saviour must
have been \emph{able} to do in order to be able to be Saviour, I grant him
unconditionally that he is right; for he speaks of a possibility; it is
a true logic of faith of which he gives a sketch. But
when he goes into particulars and wants ``by feeling'' to decide
which miraculous deeds the Saviour \emph{really} has done,
I cannot grant him that he is right, for from possibility to reality,
from the logical to the antilogical, there is a
leap.

Did Christ really walk on the sea? There is
no human being, believing any more than unbelieving, who
knows it, for of such things there is no knowledge. It is indeed
possible that Christ did not do it, although he well
could have done it. Suppose that he did not walk on the
sea, but on the other hand really stilled the storm, is he then
not the Saviour? It is certain that the same one who has
% --- p. 24 ---
\setcounter{footnote}{0}%
\opage{24}been able to raise Lazarus in Bethany\footnote[1]{Pastor Møller lays particular weight on the miracle of the raising of Lazarus, and presupposes with reason that this miracle, by my way of looking at things, must have ``full reality.'' Yes, it must have full reality! Its reality is confirmed for all believers by as strong testimonies as are well conceivable. But that in this respect it should have any peculiar advantage over so many other miracles of Jesus I cannot concede. The words: ``he stinketh,'' might easily be taken by doubters as a mere supposition of Martha's, based on Lazarus's having already ``lain four days'' in the grave. To this it must be added that some among those present went and gave the chief priests a report of such a nature that the council thereupon resolved that Jesus should die. But to want to threaten with death a man who, as is known to the eyes of all, even of the unbelieving spectators, has been able to call back to life a corpse already stinking, is not merely a sign of unbelief, but an unmistakable sign of an insanity qualified for the madhouse. One should consider, after all, what an impression such a deed must make if it had, mind you, become known to a whole assembly in an outward physical way; the report of it alone must surely strike the most hardened Sanhedrin with terror. By my way of looking at things --- it is a point I would earnestly commend to the serious reflection of believing men --- no miracle in the life of Jesus can have happened in so physically conspicuous a way and under such outwardly surveyable circumstances that no Sadducean doubter would have found a way out to a natural explanation. To perform a miracle so openly and in such a way that only eyes and ears, senses and understanding are needed, but no faith, to assure oneself of its validity as a miracle, is to abolish the religious significance of the miracle. It was precisely such an abuse of the Holy that the Devil wished to bring about when he proposed to Christ to cast himself down from the pinnacle of the Temple --- an experiment in the outwardly real, no doubt to the surprise and astonishment of a Jerusalem public.\par The negative criticism has often repeated the objection that so many and so great miraculous deeds as those of which the Evangelists tell could not possibly have really happened without Herod and Pilate and the Emperor in Rome having heard of it and having had the matter judicially investigated. This objection falls to the ground when we remember what it really means that the reality of the miracle exists only for believers. --- The grounds for this view of mine I have, with particular regard to the raising of Lazarus, set forth at greater length in the Postscript to ``The Lectures on Obstacles and Conditions for the Spiritual Life in the Present Age.''}, must also have
been able to raise the widow's son at Nain; suppose that Christ
% --- p. 25 ---
\setcounter{footnote}{0}%
\opage{25}did not raise the widow's son, but did raise Lazarus, or
the reverse: the widow's son, but not Lazarus, is he then not
the Saviour? It is a risky business, trusting in feeling,
to engage with the reasoning understanding about
the reality and the real nature of the individual
miracles in their particularity.

What, when we look closer, is no doubt Pastor
Birkedal's meaning is this, that, trusting in the credibility
of the Evangelists, one can on the basis of the impression of feeling given by the
whole judge with certainty about the particular. In the walking
on the sea, in the mastery over the storm, in the raising
of Lazarus, in all individual miracles, scenes and pictures of life,
we recognize everywhere in all the Gospels the same
figure of Christ, the same helping hand, the
same omnipotence of love: ``it is the Saviour himself, it
is like him, such a thing he has done, so he looks.''
The immediate feeling of the reality of the facts
coincides at all points with a corresponding
assurance that the Evangelists' conception of this
reality is infallibly true; but when reality
cannot fall outside the evangelical conception, hence
not outside the spirit, then it is surely indisputable
that from first to last it must be determined as Spiritual Reality.

Criticism of faith can concede changes, chronological
transpositions, expansions, abbreviations, indeed outward discrepancies
in the presentation of the particular; but in these
% --- p. 26 ---
\setcounter{footnote}{0}%
\opage{26}paraphrases there is no trace of distortion, not in the remotest way an intimation of any misunderstanding
that could cast doubt upon the validity of the conception.
The criterion of truth is that each individual picture of life evidently
corresponds to the whole basic picture: if the particular is false,
then the whole is false; if the whole is true, then the particular is
true, and that is that.

c) \emph{Reciprocal relation of criticism of knowledge and criticism of faith}.
The procedure of criticism of knowledge is correct
insofar as it aims at an approximation to objective
certainty in objective knowledge; but it cuts a ridiculous figure
by posing as if its doubt and uncertainty could
make the believer uncertain or bring about any
essential change in faith. Let it only be uncertain, from the standpoint
of objective knowledge, how it really
went with the walking on the sea; let there be, in the
name of grammar, a ``grounded supposition that''
``the episode about Peter is a later addition, owing to
the Greek translation or adaptation of the Gospel of Matthew'';
has this ``grounded supposition'' then the
slightest influence on faith? Yes, if there were an
irrefutable proof that the event is not real,
or at any rate not real in the way it is presented, that
would be another matter; then the believer, according to what has been said above
about possibility and reality, would in this particular
case alter his conception of an outwardly real thing.
But could such an alteration in the slightest degree
shake his confidence in the Evangelist's conception of the
real Christ? Does not every feature in the picture recall
the Saviour as he is? Does not ``the episode about Peter''
have something so gripping in the religious sense that every believing
% --- p. 27 ---
\setcounter{footnote}{0}%
\opage{27}soul must feel moved by it? Truly, if Christ
did not do what is told of him in this Gospel,
then he \emph{could} have done it, for --- to
speak the language of the comparison --- his other deeds in
the other Gospels correspond to it. Thus, in the worst case,
the picture of Christ as a whole would remain unchanged.
But does criticism of knowledge now have certainty that the event
cannot be real? By no means. But it does
have ``a grounded presumption.'' Yes, a presumption it
has got ``grounded'' by going around the spirit and
behind the word. Should such a presumption be able to
shake faith?

The debate between the two kinds of criticism will incessantly
turn on the relation between spirit and reality.
Criticism of knowledge strives from the account outward toward
the supernatural facts and attains the opposite of
what it seeks. It wants to seize the supernatural outside
the spirit and ends in a piecemeal dissolution of all
that is supernatural. The main edifice of the Gospels sinks into
rubble, and the fragments of reality that remain behind in all
naturalness are without significance for faith. This
must then be a spur for criticism of faith to go the
opposite way. Starting from other presuppositions, with a
different psychology, a different eye for reality and a different
measure of possibility, criticism of faith will set itself two
tasks: 1) to appropriate the scientific information
where it is in truth in its place and can count
as real information, and in addition 2) to take note
of and expose the points of view from which the errors of the understanding
proceed, points of view which, however finely the threads
may be spun in the objectively critical web, are nevertheless all
% --- p. 28 ---
\setcounter{footnote}{0}%
\opage{28}recognizable by a faith half-vitiated by knowledge and a knowledge
distorted in faith. To reason in science there corresponds
natural reality; but to the spirit of revelation there corresponds
Spiritual Reality.

\begin{center}V.\end{center}

\emph{Spiritual Reality is a conjunction of
extremes, whose middle affords the possibility, not
merely of a reconciliation between faith and knowledge}, but
\emph{also of a separation of faith and opinion.}

The reality of revelation, whether we look at the fact
or at the apprehension, is Spiritual Reality; but
the point of view of natural reality is hereby by no means
excluded. This point of view will therefore also assert its right
in the sacred history. Concerning the Creation, Paradise,
the Fall and the End of the World, criticism of knowledge is
entitled to claim for natural reality a significance
which it has not received in the biblical presentation.
Since reality stands against reality (natural reality
against that presented in the word of revelation),
knowledge retains its right over against faith; but the
conflict thereby arising is resolved in the cognition that
the dualism is posited by the spirit itself. The same almighty
spirit of God that has created the world has also imparted
the word of the creation of the world. That God's unconditionally free
will is the cause of this world is the truth which
the spirit, by its word of revelation, wishes to impress upon the human race
in a manner comprehensible and intuitive to all; the infinite
reciprocal action of the forces of the world, the essence of the laws
and the interconnection of natural things, on the other hand, were not to be
an object of revelation; the investigations belonging
% --- p. 29 ---
\setcounter{footnote}{0}%
\opage{29}here are left to reason and to science. Criticism of knowledge
is therefore right when it points out
that creation in six days does not agree with the course of natural development
which can in reality be demonstrated.
But it is wrong when it seeks the ground of the discrepancy
in the fact that the spirit which speaks through
the history of creation has not been versed in modern physics
and has therefore shown itself ignorant of the real
interconnection of things. Criticism of faith is right when it proceeds
from the presupposition that the spirit which expresses itself in
the word of the Creation is the almighty, omniscient spirit,
the power of the Father that has brought forth heaven and earth;
but it is wrong when, by borrowing a little from
natural science and tinkering with the revealed word,
it wants the history of creation adapted to astronomy.\footnote[1]{Compare, Philosophy of Religion, pp.\ 151---66.} % print: Jevnf, (as printed)
It is the theological muddling, the imagined unity
of faith and knowledge, that forms the point of departure for both
errors.

At the stage of revelation the activity of the spirit
by which the word in the history of creation was imparted is an act
as divine as that by which the world was
created. The extremes signify that the act of creation is in itself
incomprehensible to man, that therefore the spirit through the word must
enlighten the believers, must meet their capacity of comprehension and religious
need. The God who has revealed his omnipotence in nature has also revealed his will in
the word: revelation itself is the middle of the extremes.

As it is with the history of creation, so
it is with every section of the sacred
% --- p. 30 ---
\setcounter{footnote}{0}%
\opage{30}history where the reality expressed in the physical order of the world
stands in opposition to the revealed word:
here is an essential opposition, but no real conflict.
In the apprehension and judgment of the middle corresponding to the extremes
there now opens a possibility for the separation
between \emph{faith and opinion}, a separation which hitherto,
unfortunately, has been overlooked in the doctrine of religion all too much, not to say entirely.

\emph{Faith in the Word in} the history of creation is and must
be the same, unalterably the same, whether the
believer otherwise belongs to those who know or to those who do not know.
But the individual apprehension of the details of the presentation
is not the same. He who does not know takes
everything immediately according to the letter, without any inkling that here
there either should or can be any thought of a separation between
natural reality in the outward facts and Spiritual Reality
in the word. He who knows, on the other hand, has made
the opposition clear to himself; but he has also made it
clear to himself that in the biblical presentation, simply understood,
there is not a jot or tittle that in any way
could weaken faith in the divinity of the word. So
they both, he who knows and he who does not know,
have the same faith, but it is far from the case that they have the same
opinion.

The categorical difference between faith and opinion is
easy to demonstrate. To take yet a couple of examples
along, we choose Paradise and the Fall.

So much idle strife has now been conducted about the
geographical location of Paradise, the nature of
the tree of knowledge and the diabolical serpent
that it might well soon be enough. The endless squabbling falls away
% --- p. 31 ---
\setcounter{footnote}{0}%
\opage{31}when one reflects on the difference between faith
and opinion. To be at one in faith, that is: to be
at one in the religious appropriation of the word, that is: to be
completely at one in the apprehension of the co-created image of God,
the ideal of innocence, the temptation, the Fall, the guilt and
the punishment, an apprehension which can and must be exactly determined
by the individual features and articulation of the narrative. It
is the first decisive presuppositions of religious psychology
that are at issue. If there is agreement on this
point, then there is agreement in faith. What outward reality
may be attributed to the trees, and how the serpent
looked, on that opinions may be divided. The
individual representations that accompany the impression of the
whole are so far conditioned by the feeling,
imagination and stage of cultivation of each individual that the one dare not
force his opinions upon the other.

How important it is everywhere in the domain of religious views
to make a thorough distinction between faith and
opinion will also be evident from the many years' long
polemic for and against the Grundtvigian view of the
origin of the creed and its relation to the baptismal word,
which runs like a red thread through our most recent
church history. ``The vital nerve of the Grundtvigian view''
must, from my point of departure, be sought expressly in this, ``that
faith in the creed in principle coincides
with faith in its origin from the object of faith,
that is, with faith in the Christ personally
present in the words of the confession. (Philosophy of Religion, p.\ 72). Whatever particular
representations ``Grundtvigians'' and ``anti-Grundtvigians''
may wish to set in circulation concerning the outward
reality in the manner in which Christ in
% --- p. 32 ---
\setcounter{footnote}{0}%
\opage{32}his ``days of flesh'' is supposed to have imparted or not imparted to
the Apostles the wording of the creed, I must still
repeat my simplicity: there is no human being on
earth who knows that; for of such things there is no objective
knowledge. The task then becomes to investigate what
in this matter has decisive significance for faith, and
what may be regarded as opinions and be either
reasonable or unreasonable conjectures. That Grundtvig has not
overlooked the difference between faith and opinion is seen,
among other things, from the following statement: ``We who call ourselves
Christ's messengers, and who are to impart to the congregation all
the enlightenment in faith which the spirit by Scripture or
in some other way bestows on us, must well beware of
making our enlightenments, whether about the creed as a
whole or about its individual members and parts, into articles of faith.''\footnote[1]{Den kristelige Børnelærdom [Christian Instruction of Children] 1868. p.\ 45.}

Which determinations, then, can be fixed as belonging
to faith? They can only be such as
indisputably belong to the ground of faith and furnish the first
indispensable conditions for the life of faith.

The baptismal word is a sacramental word. Faith in the sacramental word
includes the presupposition that not only is baptism performed
according to Christ's command, but that the act of baptism itself in
the spirit through the word is an act of Christ.\footnote[2]{Compare Philosophy of Religion, pp.\ 508---514.} The words:
I baptize thee in the name of the Father and of the Son and of the
Holy Ghost, must then be apprehended as Christ's words, words ``from
the Lord's own mouth'' to him who lets himself be baptized. But
he who says: ``I baptize thee''\dots{}. must surely be the
% --- p. 33 ---
\setcounter{footnote}{0}%
\opage{33}same as he who beforehand has set the conditions of baptism
and received the promise of the fulfilment of these conditions through
a sincere and voluntary answering of the questions:
``Dost thou renounce\dots{}.? Dost thou believe\dots{}.? Dost thou believe\dots{}.?
Dost thou believe\dots{}.? To put such words into Christ's mouth,
and yet secretly to harbor the subsidiary notion that they
are not really his own words, but words of another origin,
words that with a certain well-meaning arbitrariness are ascribed
to him, in order to give the act greater solemnity, what
is that? It is to reduce the sacrament to a ceremony,
it is to hollow out the basic word and to begin the life of faith with
a self-contradiction in faith.

The apprehension of the sacrament as an act of Christ
in the congregation follows with irrefutable consistency
from the principle of faith, and Pastor C. Møller hits exactly
upon my way of looking at it when he reminds us that nothing can be
fixed about faith except what can be derived from the principle.

But now the divine basic word of the creed
surely also has a human side. About the sentences in the
three apostolic articles, as regards manner of transmission,
expressions of translation and the like, philological
and historical-critical investigations can of course be undertaken, and
about these again different opinions can form. To what extent
these opinions are really subordinate to faith will
depend on whether from all sides the aim is at
the original. In the vigilant and persistent effort
to secure and preserve the form of the creed in
every individual member and expression, corresponding to the original
word of faith, the creative word of life, ``the word from the Lord's
own mouth,'' we should then see not a one-sided zeal
for a party cause, but a sincere love of truth.

% --- p. 34 ---
\setcounter{footnote}{0}%
\opage{34}To make the creed into a kind of incomplete
compendium of a ``Christian dogmatics'' and to take the content
so generally that one can more or less at one's pleasure
change members and expressions, is to go to an
extreme; to want, by a brand-new revelation, to secure
a literal reproduction of the linguistic form, and
the outward manner in which the basic word was originally
imparted, is to go to an opposite extreme. In the
first extreme one goes around the word in order to grasp
the spirit; in the last extreme one goes further than the spirit
in order to prop up the word. The truth is the middle
of the extremes, is Spiritual Reality, reality for the spirit.

\begin{center}VI.\end{center}

\emph{Spiritual Reality, the mark of truth for
the creed and Holy Scripture, makes
historical credibility absolutely dependent
on religious credibility and gives criticism of faith a measure
for distinguishing revelation from mythology}.

However great the distance between the seventeenth century's
scriptural-dogmatic doctrine of revelation and the present day's
Bible-critical mythology, there is nevertheless one principle
that is common to both.

``Every inner alteration, conditioned by appropriation,
which an outward event undergoes in the narrator's
presentation, \emph{is a falsification} of its reality'';
this is the supreme premise, the logical major premise,
from which orthodox apologetics and
negative criticism have each formulated their conclusion. ---
``The sacred writers have not let the events
% --- p. 35 ---
\setcounter{footnote}{0}%
\opage{35}undergo any alteration whatsoever; \emph{ergo}: the sacred
history is in all points and parts a history in literal agreement
with itself, outwardly real.'' Thus
orthodoxy concludes. --- ``The biblical authors have
in every presentation of the supernatural let the event
undergo a transformation, if unconsciously, yet essential;
\emph{ergo}: the whole so-called sacred history
is in everything essential a falsified, unreal history.''
Thus negative criticism concludes. Mediating theology
mediates and mediates, but without bearing, without principle.

Is it then really the case that every alteration brought about
by believing appropriation, which a supernatural event can
undergo, is a falsification? This fundamental question, decisive for the credibility of the Bible, must be answered
with Yes or No.

If the question is answered with Yes, then the biblical word is
mythical, and dissolving criticism of knowledge is in the main
right. ``The possibility of myth becomes probability
when the wondrous in the narratives, the irresolvable
contradictions between the individual Evangelists, explicable only
by the most artificial harmonistics, the manifold
chronological difficulties that cannot be
brought into agreement with contemporary secular writers,
are taken into consideration. Nor is it
sufficient, as has hitherto been the practice, to settle all these
contradictions by abandoning one Evangelist and favoring
another, so that Matthew is sacrificed for
Luke, and the latter, when required, again for John.
This is merely to measure with unequal measure and weight. It
moreover appears, on an impartial consideration of the different
historical presentations, that all the Evangelists must
% --- p. 36 ---
\setcounter{footnote}{0}%
\opage{36}be disallowed in equal degree, that the testimony of the one is
worth just as much, or rather, just as little as
that of the other. Here a bellum omnium contra omnes is waged,
and nowhere are we able to win firm historical ground.''\footnote[1]{K. Schwarz. The History of the Most Recent Theology in Brief Presentation (Den nyeste Theologies Historie i kort Fremstilling; translated from the third edition of the original, Copenhagen 1869). p.\ 94.}

If the question is answered with No, then the biblical word enters
anew into force as the word of revelation; the doctrine of faith
has by this No freed itself from the objections and imputations of dissolution-criticism
and of mythical interpretation.
It lies in the nature of the matter, as shown above, that miracle
and faith must be inseparable in apprehension and presentation;
the transformation that the apprehension has brought about is
a transformation in the spirit. The transformation is a
falsification only when the spirit is false. The mythical spirit,
which posits the unreal as real, is a false spirit.
But the spirit of revelation is the spirit of truth, because
in the deepest sense it is the spirit of reality.

That, as regards the outward and the particular, there can
be mistakes which do not in the remotest degree falsify
the content of faith, only a spiritless, letter-bound supranaturalism
can deny. From the point of view of the faith of the Gospel
a principled distinction can then also be made between legend,
legendary features, and myths. How ``legends could get
out among the brethren,'' and even legends that were
occasioned by an innocent misunderstanding of Jesus' words,
is seen from John. 21:23. But such legends do not intermingle
any element deviating from the thought of revelation and conflicting with its essence,
do not falsify the principle of reality
and are consequently not myths. In Matth. 27:51 it is told,
% --- p. 37 ---
\setcounter{footnote}{0}%
\opage{37}that when Jesus ``gave up the spirit on the cross, the earth quaked
and the rocks were rent, and the graves were opened, and many
bodies of the saints that slept arose, and went out
of the graves after his resurrection, and came into the
holy city and appeared unto many.'' This can,
if one really wishes, be apprehended as a legend, but it can
in no way be regarded as a myth. If the resurrection
of the dead of Jerusalem is a myth, then the principle of revelation is
denied, and then the raising of Lazarus and
the resurrection of Christ are also myths, for the one proceeds
from the same principle as the other.\footnote[1]{That the relation, however, cannot be reversed, so that Christ's resurrection became a legend, and only the resurrection of the bodies of the saints that slept a reality, is understood of itself.}

It is an easy matter to make the biblical narratives
count as full reality when the whole thing is accomplished
by taking everything without further ado according to the letter,
leaning on blind authority and putting assertions in place of
grounds. What is easier than to say: ``It stands in the
Bible, it is historical, and \emph{therefore} we are to believe it;
unless we had beforehand irrefutable proofs that
Holy Scripture is more exact than any legal document,
our faith would have no secure foundation; but it
stands written, \emph{it stands written, it stands literally
written}! --- By opening the door of contradiction and
lighting the juridical assertiveness of Scripture down the stairs,
criticism of knowledge, humanly speaking, has done the biblical word a true
service.

There is reality enough, ``full reality'' in the revealed
facts; the sacred history is like a
richly laid table: step in, for all is ready. But
% --- p. 38 ---
\setcounter{footnote}{0}%
\opage{38}what is thus made ready benefits only those who worthily
understand how to enjoy it. To declare the supernatural nourishments
to be dishes for show, because one, by inhaling
poison, has lost the power of digestion, is tragically frivolous, wanton.
To sit down at the table and fill one's dogmatic pockets
with outward matter which one intends to stuff into others,
but cannot digest oneself, is a crude misuse of what is
holy. The reality of revelation is reality only
for those who in simplicity of heart have the need of appropriation
and in the understanding of faith the gift of appropriation. Of other
reality one may, as the case may be, soon get
too little, soon again too much, unfortunately all too much;
but of the reality of revelation one can in faith always
be sure of getting neither too little nor too much;
for appropriation is the measure, and appropriation is conditioned by
need. The reality of revelation is, in the outward as
in the inward, Spiritual Reality, a reality which, on the conditions of
appropriation and self-understanding, has been from the very beginning,
is, and in all eternity will be reality
for the spirit.
\signature{R. \emph{Nielsen}.}

\end{document}
